%=======================================================================
% GIST Master's Thesis Template (English) - main file
%
%  * Compiler : pdfLaTeX  (Overleaf: Menu > Compiler > pdfLaTeX)
%  * Bibliography : BibTeX (references.bib, style ieeetr)
%  * 수정할 곳은 [EDIT] 표시를 검색하면 모두 찾을 수 있습니다.
%  * 자세한 사용법은 README.pdf를 참고하세요.
%=======================================================================
\documentclass{gist}

%-----------------------------------------------------------------------
% Packages
%-----------------------------------------------------------------------
\usepackage{type1cm}          % 임의 글자 크기(\fontsize) 지원 - 표지에 필요
\usepackage{graphicx}
\usepackage{subcaption}       % (a), (b) 부분 그림
\usepackage{rotating}         % sidewaysfigure / sidewaystable (가로 페이지)
\usepackage{placeins}         % \FloatBarrier
\usepackage{float}            % [H] 위치 옵션
\usepackage{caption}

\usepackage{amsmath, amssymb}

\usepackage{array}
\usepackage{booktabs}         % \toprule, \midrule, \bottomrule
\usepackage{multirow}
\usepackage{makecell}
\usepackage{tabularx}
\usepackage{tabularray}       % tblr / longtblr 표
\UseTblrLibrary{booktabs}

\usepackage{enumitem}
\usepackage{url}
\usepackage{setspace}
\usepackage{appendix}

% [EDIT/확인] 아래 geometry 설정은 gist.cls에 정의된 여백(본문 폭 155 mm 등)을
% 덮어씁니다(사방 1 inch). 이 상태로 인준을 받은 선례가 있지만, 제출 전에
% 학교의 최신 학위논문 작성 지침과 대조해서 유지 또는 삭제를 결정하세요.
\usepackage[margin=1in]{geometry}

\usepackage[square, numbers, sort&compress]{natbib}

\usepackage[hidelinks, hypertexnames=false]{hyperref}  % hypertexnames=false: 표지·초록·본문 쪽 번호가 겹쳐도 링크가 꼬이지 않게 함
\usepackage{cleveref}         % \Cref{...} : "Fig. 3.1", "Table 4.2" 자동 생성
\crefname{figure}{fig.}{figs.}
\crefname{subfigure}{fig.}{figs.}
\crefname{equation}{eq.}{eqs.}
\Crefname{figure}{Fig.}{Figs.}
\Crefname{subfigure}{Fig.}{Figs.}
\Crefname{equation}{Eq.}{Eqs.}

\setcounter{secnumdepth}{3}

% 자주 쓰는 약어 매크로
\newcommand\etal{\textit{et~al.\ }}
\newcommand\eg{\textit{e.g.\ }}
\newcommand\ie{\textrm{i.e.\ }}

%-----------------------------------------------------------------------
% (선택) 필요할 때만 주석을 해제하세요.
%-----------------------------------------------------------------------
% 알고리즘 의사코드 - 사용 시 아래 \listalgorithms 주석도 해제
%\usepackage[ruled,vlined,linesnumbered]{algorithm2e}
%\usepackage{etoolbox}
%\AtBeginEnvironment{algorithm}{\setstretch{1.0}}

% 정리(Theorem) 환경
%\usepackage{theorem}
%\newtheorem{theorem}{Theorem}[section]
%\newtheorem{lemma}[theorem]{Lemma}
%\newtheorem{corollary}[theorem]{Corollary}

% 기타: 대각선 표 머리(diagbox), 옆 캡션(sidecap), 가로 페이지 전체(pdflscape)
%\usepackage{diagbox}
%\usepackage[rightcaption]{sidecap}
%\usepackage{pdflscape}

%-----------------------------------------------------------------------
% References를 목차(Contents)에 표시
% natbib이 gist.cls의 참고문헌 정의를 덮어써서 목차 항목이 빠지는 문제를 보완합니다.
%-----------------------------------------------------------------------
\renewcommand{\bibsection}{%
  \chapter*{\bibname}%
  \phantomsection
  \addcontentsline{toc}{chapter}{\fontsize{12}{1pt}\selectfont \bibname}}

%=======================================================================
% 논문 정보 [EDIT]
%=======================================================================

%-----------------------------------------------------------------------
% Department code [EDIT]
% 형식: \code{{학위/}{학과코드}}   학위: BS/ (학사), MS/ (석사), PhD/ (박사)
% 학과 코드별로 표지·초록에 찍히는 학과명 (gist.cls에 정의된 그대로)
% IC - School of Information and Communications (정보통신공학부)
% IM - School of Information and Communications (정보기전공학부)
% EC - School of Electrical Engineering and Computer Science (전기전자컴퓨터공학부)
% MS - School of Materials Science and Engineering (신소재공학부)
% ME - School of Mechanical and Robotics Engineering (기계로봇공학부)
% EN - Department of Environment and Energy Engineering (환경에너지공학과)
% LS - School of Life Sciences (생명과학부)
% PH - Department of Physics and Photon Science (물리·광과학과)
% CH - Department of Chemistry (화학과)
% NA - Department of Nanobio Materials and Electronics (나노바이오재료전자공학과)
% MD - Department of Biomedical Science and Engineering (의생명공학과)
% ET - School of Integrated Technology (에너지융합학제전공)
% CT - School of Integrated Technology (문화기술융합학제전공)
% RT - School of Integrated Technology (지능로봇프로그램전공)
% 일부 학과명은 현행 명칭과 다를 수 있으니, 컴파일 후 표지의 학과명을 꼭 확인하세요.
% 학과명이 바뀌었거나 목록에 없으면 gist.cls의 \ifx\@dept\EN ... \fi 블록을 수정하세요.
%-----------------------------------------------------------------------
\code{{MS/}{EN}}

%-----------------------------------------------------------------------
% Thesis title in English [EDIT]
% 줄을 나누고 싶은 위치에 \titlebreak를 넣으세요. '\\'는 쓰지 마세요.
\etitle{Title of Your Master's Thesis in English}

% Thesis title in Korean [EDIT]
\ktitle{국문 학위논문 제목}

%-----------------------------------------------------------------------
% Advisor's name, without a position such as 'Prof.' [EDIT]
\advisor{Advisor Name}
\kadvisor{지도교수}

% Co-advisor (없으면 주석 처리 유지)
%\coadvisor{Co-advisor Name}
%\kcoadvisor{공동지도교수}

%-----------------------------------------------------------------------
% Name of the author [EDIT]
\ename{Gildong Hong}

% 한글 이름: 7칸 고정, 한 글자씩 {}로 감싸고 오른쪽 끝에 맞춰 입력합니다.
% 예) 세 글자 이름 -> {}{}{}{}{홍}{길}{동},  네 글자 이름 -> {}{}{}{남}{궁}{길}{동}
\kname{{}{}{}{}{홍}{길}{동}}

%-----------------------------------------------------------------------
% Student ID [EDIT]
\studentid{2024XXXX}

%-----------------------------------------------------------------------
% The year of graduation (ex. 2027) [EDIT]
\coveryear{YYYY}

% The date signed by the advisor: {Month}{Day}{Year}  ex) {December}{10}{2026} [EDIT]
\advisorsigndate{Month}{DD}{YYYY}

% The date signed by the referees (최종 서명일) [EDIT]
\refereesigndate{Month}{DD}{YYYY}

%-----------------------------------------------------------------------
% Names of the referees [EDIT]
% 석사: refereeA ~ refereeC (3인), 박사: refereeA ~ refereeE (5인)
% refereeA는 Committee Chair로 표시되며, 보통 지도교수가 맡습니다.
\refereeA{Prof. Committee Chair}
\refereeB{Prof. Committee Member 1}
\refereeC{Prof. Committee Member 2}
%\refereeD{Prof. Committee Member 3}   % 박사
%\refereeE{Prof. Committee Member 4}   % 박사

%=======================================================================
\begin{document}
%=======================================================================

%-----------------------------------------------------------------------
% Abstract (English) -> misc/eabstract.tex
%-----------------------------------------------------------------------
\begin{eabstract}
\begin{spacing}{2.0}
%%%%%%%%%%%%%%%%%%%%%%%%%%%%%%%%%%%%%%
% Abstract (English)  [EDIT]
% - 첫 줄의 학위/학과 코드, 학번, 이름, 제목, 쪽수, 지도교수는 thesis.tex의 정보로 자동 생성됩니다.
% - 여기에는 초록 본문만 작성하세요. 빈 줄로 문단을 나눕니다.
%%%%%%%%%%%%%%%%%%%%%%%%%%%%%%%%%%%%%%
Write the background and motivation of your research here. Briefly state the problem that this thesis addresses and why it matters.

Describe the data and methods you used, and summarize the main results with key quantitative findings where possible.

Conclude with the significance of the results and their implications for the field.

\end{spacing}
\end{eabstract}

%-----------------------------------------------------------------------
% (선택) 국문요약 -> misc/kabstract.tex
% 영문 논문에 국문요약이 필요한지는 학과·학교 지침을 확인하세요.
% 필요하면 아래 5줄의 주석을 해제하면 됩니다.
%-----------------------------------------------------------------------
%\begin{kabstract}
%\begin{spacing}{2.0}
%%%%%%%%%%%%%%%%%%%%%%%%%%%%%%%%%%%%%%%
% 국문요약 (선택)  [EDIT]
% thesis.tex에서 kabstract 블록의 주석을 해제해야 출력됩니다.
%%%%%%%%%%%%%%%%%%%%%%%%%%%%%%%%%%%%%%
연구의 배경과 목적을 기술합니다.

사용한 자료와 방법, 주요 결과를 요약합니다.

연구의 의의와 활용 방안을 기술합니다.

%\end{spacing}
%\end{kabstract}

%-----------------------------------------------------------------------
% Contents, list of tables, list of figures
%-----------------------------------------------------------------------
\makecontents
\listtables      % 표가 없으면 주석 처리
\listfigures     % 그림이 없으면 주석 처리
%\listalgorithms % algorithm2e 사용 시 주석 해제

%=======================================================================
% Main body
% 각 장은 chapters/ 폴더의 파일로 분리되어 있습니다.
% 장을 추가·삭제하려면 아래 블록을 복사하거나 지우세요.
%=======================================================================

% DO NOT DELETE THE FOLLOWING TWO LINES (본문 쪽 번호를 1부터 시작)
\begin{spacing}{2.0}
\pagenumbering{arabic}
\setcounter{page}{1}

%%%%%%%%%%%%%%%%%%%%%%%%%%%%%%%%
% Chapter 1. Introduction
%%%%%%%%%%%%%%%%%%%%%%%%%%%%%%%%
\chapter{Introduction} \label{chapter1}

%==============================%
% 1.1 Research background
%==============================%
\section{Research background} \label{chap1:sec1}

Describe the broad context of your research here. Every chapter, section, figure, table, and equation should have a \verb|\label{...}| so that it can be referenced automatically.

To cite a reference, use the \verb|\cite| command with an entry key from the bibliography file (\texttt{references.bib}). A single citation looks like this \cite{RN35}, and multiple citations are sorted and compressed automatically \cite{haralick1973a, akiba2019optuna, zhou_unet_2018}. Online sources such as databases can also be cited \cite{nier_weis}.

%==============================%
% 1.2 Problem statement
%==============================%
\section{Problem statement} \label{chap1:sec2}

State the limitations of existing approaches and the specific gap that this thesis addresses. Commonly used abbreviation macros are available, \eg \verb|\eg|, \verb|\ie|, and \verb|\etal| (\eg Haralick \etal \cite{haralick1973a}).

%==============================%
% 1.3 Research objectives
%==============================%
\section{Research objectives} \label{chap1:sec3}

List the objectives of the thesis. Use \texttt{itemize} or \texttt{enumerate} for lists:
\begin{enumerate}[label=(\arabic*)]
    \item The first objective of this thesis.
    \item The second objective of this thesis.
\end{enumerate}

%==============================%
% 1.4 Thesis organization
%==============================%
\section{Thesis organization} \label{chap1:sec4}

Cross-references are written with \verb|\Cref{label}|, which inserts the type and number automatically. The rest of this thesis is organized as follows. \Cref{chapter2} reviews related work. \Cref{chapter3} describes the data and methods. \Cref{chapter4} presents and discusses the results. Finally, \Cref{chapter5} concludes the thesis.
   % Chapter 1: Introduction
\end{spacing}

\begin{spacing}{2.0}
%%%%%%%%%%%%%%%%%%%%%%%%%%%%%%%%
% Chapter 2. Literature Review
%%%%%%%%%%%%%%%%%%%%%%%%%%%%%%%%
\chapter{Literature Review} \label{chapter2}

%==============================%
% 2.1 Topic A
%==============================%
\section{Review topic A} \label{chap2:sec1}

Summarize previous studies on the first topic related to your research \cite{RN35}.

%==============================%
% 2.2 Topic B
%==============================%
\section{Review topic B} \label{chap2:sec2}

Summarize previous studies on the second topic \cite{zhou_unet_2018}.

%==============================%
% 2.3 Research gaps
%==============================%
\section{Research gaps and contributions} \label{chap2:sec3}

Identify the research gaps from the review above and explain how this thesis contributes to filling them.
   % Chapter 2: Literature Review
\end{spacing}

\begin{spacing}{2.0}
%%%%%%%%%%%%%%%%%%%%%%%%%%%%%%%%
% Chapter 3. Data and Methods
% 예시: 그림(부분 그림 라벨), 가로 페이지 그림, 표(tblr), 수식
%%%%%%%%%%%%%%%%%%%%%%%%%%%%%%%%
\chapter{Data and Methods} \label{chapter3}

%==============================%
% 3.1 Study area
%==============================%
\section{Study area} \label{chap3:sec1}

% ------------------------------------------------------------------
% [예시 1] 여러 패널이 하나의 이미지 파일에 들어 있는 그림
% 패널 (a), (b)를 개별적으로 참조하려면 \phantomsubcaption으로 라벨만 만들어 둡니다.
% 그림 파일은 figures/ 폴더에 넣고 \includegraphics{figures/파일명.pdf}로 불러오세요.
% 아래 example-image-a는 TeX에 내장된 예시 이미지이므로, 실제 그림으로 바꾸면 됩니다.
% \caption[짧은 제목]{긴 캡션}: 대괄호 안의 짧은 제목만 그림 목차(List of Figures)에 들어갑니다.
% ------------------------------------------------------------------
\begin{figure}[!htbp]
    \centering
    {
      \phantomsubcaption\label{fig:study_area_a}
      \phantomsubcaption\label{fig:study_area_b}
    }
    \includegraphics[width=0.8\textwidth]{example-image-a}
    \caption[Overview of the study area.]{Overview of the study area. (a) Location of the study area and (b) land cover map. Only the short title given in the square brackets of the caption command appears in the List of Figures.}
    \label{fig:study_area}
\end{figure}

Describe your study area here. \Cref{fig:study_area} shows the study area, and individual panels can be referenced as \Cref{fig:study_area_a} or \Cref{fig:study_area_a,fig:study_area_b}.

%==============================%
% 3.2 Datasets
%==============================%
\section{Datasets} \label{chap3:sec2}

Describe the datasets. \Cref{tab:dataset} summarizes the datasets used in this study.

%----------------------------------------------------------
% [표 예시 1] tabularray(tblr): 열 너비 자동 조정, 세로 병합
%  - colspec: Q = 내용 폭, X = 남은 폭을 채움 / [c,m] = 가운데·세로 중앙
%  - cell{2}{3} = {r=2}{} : 2행 3열 셀을 아래로 2칸 병합(병합되는 칸은 비워 둠)
%----------------------------------------------------------
\begin{table}[!htbp]
\centering
\small
\caption{Summary of datasets (example of a \texttt{tblr} table).}
\label{tab:dataset}
\begin{tblr}{
  width = \linewidth,
  colspec = { Q[c,m] X[l,m] Q[c,m] Q[c,m] },
  row{1} = {font=\bfseries, halign=c},
  cell{2}{3} = {r=2}{},
}
\toprule
Dataset   & Description                                   & Type          & Resolution \\
\midrule
Dataset A & Description of dataset A                      & {Satellite \\ imagery} & 10 m \\
\addlinespace
Dataset B & Description of dataset B                      &               & 30 m \\
\addlinespace
Dataset C & Description of dataset C (e.g., in-situ data) & Point data    & -- \\
\bottomrule
\end{tblr}
\end{table}


%==============================%
% 3.3 Methodology
%==============================%
\section{Methodology} \label{chap3:sec3}

\subsection{Overall workflow} \label{chap3:sec3:subsec1}

The overall workflow is shown in \Cref{fig:workflow}. A wide flowchart can be placed on a rotated page with \texttt{sidewaysfigure}.

% ------------------------------------------------------------------
% [예시 2] 가로로 회전한 페이지 전체를 차지하는 그림 (rotating 패키지)
% ------------------------------------------------------------------
\begin{sidewaysfigure}
    \centering
    \includegraphics[width=0.8\textheight]{example-image}
    \caption{Flowchart of the methodology (example of a rotated, full-page figure).}
    \label{fig:workflow}
\end{sidewaysfigure}

\subsection{Model description} \label{chap3:sec3:subsec2}

\subsubsection{Equations} \label{chap3:sec3:subsec2:subsubsec1}

% ------------------------------------------------------------------
% [예시 3] 번호가 붙는 수식. 본문에서는 \Cref{eq:...}로 참조합니다.
% ------------------------------------------------------------------
Numbered equations are written with the \texttt{equation} environment and referenced as \Cref{eq:rmse}:

\begin{equation}
\label{eq:rmse}
RMSE = \sqrt{\frac{1}{n} \sum_{i=1}^{n} \left(y_i - \hat{y}_i\right)^2}
\end{equation}

where $n$ is the number of samples, $y_i$ is the observed value, and $\hat{y}_i$ is the predicted value. Piecewise definitions use \texttt{cases}, as in \Cref{eq:piecewise}:

\begin{equation}
\label{eq:piecewise}
f(x) =
\begin{cases}
    1, & \text{if } x \geq \tau \\
    0, & \text{otherwise}
\end{cases}
\end{equation}

Multi-line equations aligned at the equal sign use \texttt{align} (\Cref{eq:precision,eq:recall}):

\begin{align}
Precision &= \frac{TP}{TP+FP} \label{eq:precision} \\
Recall    &= \frac{TP}{TP+FN} \label{eq:recall}
\end{align}

\subsubsection{Evaluation metrics} \label{chap3:sec3:subsec2:subsubsec2}

Describe the evaluation metrics and validation strategy here. Detailed settings can be moved to the appendix (see \Cref{tab:hyperparameters}).
   % Chapter 3: Data and Methods
\end{spacing}

\begin{spacing}{2.0}
%%%%%%%%%%%%%%%%%%%%%%%%%%%%%%%%
% Chapter 4. Results and Discussion
% 예시: 표(booktabs + multirow), subfigure 환경을 이용한 부분 그림
%%%%%%%%%%%%%%%%%%%%%%%%%%%%%%%%
\chapter{Results and Discussion} \label{chapter4}

%==============================%
% 4.1 Results
%==============================%
\section{Results} \label{chap4:sec1}

\subsection{Quantitative results} \label{chap4:sec1:subsec1}

\Cref{tab:results} compares the performance of the models. Tables with grouped rows and columns can be built with \texttt{booktabs}, \texttt{multirow}, and \textbackslash\texttt{cmidrule}.

%----------------------------------------------------------
% [표 예시 2] tabular + booktabs + multirow
%  - \multirow{2}{*}{...} : 세로 2칸 병합
%  - \multicolumn{2}{c}{...} : 가로 2칸 병합
%  - \cmidrule(lr){3-4} : 3~4열에만 가로선
%----------------------------------------------------------
\begin{table}[!htbp]
\centering
\small
\caption{Performance comparison (example of a \texttt{booktabs} table).}
\label{tab:results}
\setlength{\tabcolsep}{6pt}
\begin{tabular}{cc cc cc}
\toprule
\multirow{2}{*}{\textbf{Model}} & \multirow{2}{*}{\textbf{Input}} &
\multicolumn{2}{c}{\textbf{Site A}} &
\multicolumn{2}{c}{\textbf{Site B}} \\
\cmidrule(lr){3-4} \cmidrule(lr){5-6}
& & Metric 1 & Metric 2 & Metric 1 & Metric 2 \\
\midrule
\multirow{2}{*}{Model 1}
& Baseline      & 0.00 & 0.00 & 0.00 & 0.00 \\
& \textbf{Proposed} & 0.00 & 0.00 & \textbf{0.00} & \textbf{0.00} \\
\addlinespace
\multirow{2}{*}{Model 2}
& Baseline      & 0.00 & 0.00 & 0.00 & 0.00 \\
& \textbf{Proposed} & 0.00 & 0.00 & \textbf{0.00} & \textbf{0.00} \\
\bottomrule
\end{tabular}
\end{table}


\subsection{Qualitative results} \label{chap4:sec1:subsec2}

% ------------------------------------------------------------------
% [예시 4] 그림 파일이 패널별로 따로 있을 때: subfigure 환경
% 각 패널 아래에 (a), (b) 캡션이 자동으로 붙습니다.
% ------------------------------------------------------------------
\begin{figure}[!htbp]
    \centering
    \begin{subfigure}[b]{0.48\textwidth}
        \centering
        \includegraphics[width=\textwidth]{example-image-b}
        \caption{Baseline}
        \label{fig:map_baseline}
    \end{subfigure}
    \hfill
    \begin{subfigure}[b]{0.48\textwidth}
        \centering
        \includegraphics[width=\textwidth]{example-image-c}
        \caption{Proposed}
        \label{fig:map_proposed}
    \end{subfigure}
    \caption{Visual comparison of results: (a) baseline and (b) proposed method.}
    \label{fig:maps}
\end{figure}

\Cref{fig:maps} compares the results visually. The proposed method (\Cref{fig:map_proposed}) shows ... compared with the baseline (\Cref{fig:map_baseline}).

%==============================%
% 4.2 Discussion
%==============================%
\section{Discussion} \label{chap4:sec2}

Interpret the results, compare them with previous studies, and discuss their implications.

\subsection{Limitations and future work} \label{chap4:sec2:subsec1}

Describe the limitations of this study and directions for future research.
   % Chapter 4: Results and Discussion
\end{spacing}

\begin{spacing}{2.0}
%%%%%%%%%%%%%%%%%%%%%%%%%%%%%%%%
% Chapter 5. Conclusion
%%%%%%%%%%%%%%%%%%%%%%%%%%%%%%%%
\chapter{Conclusion} \label{chapter5}

Summarize the main findings of the thesis, restate the contributions, and suggest directions for future research.

%-----------------------------------------------------------------------
% 참고: 연구가 독립된 두 주제(Part I, Part II)로 구성된 경우,
% 다음과 같이 장 제목을 구성하는 방법도 있습니다.
%   Chapter 1  Introduction
%   Chapter 2  Literature Review
%   Chapter 3  Part I: [첫 번째 연구 주제]   (Data and methods / Results / Discussion)
%   Chapter 4  Part II: [두 번째 연구 주제]  (Data and methods / Results / Discussion)
%   Chapter 5  Overall Conclusion
%-----------------------------------------------------------------------
   % Chapter 5: Conclusion
\end{spacing}

%-----------------------------------------------------------------------
% (선택) Summary 페이지 -> misc/summary.tex
%-----------------------------------------------------------------------
%\begin{spacing}{2.0}
%%%%%%%%%%%%%%%%%%%%%%%%%%%%%%%%%%%%%%%
% Summary (선택)
% thesis.tex에서 summary 블록의 주석을 해제해야 출력됩니다.
%%%%%%%%%%%%%%%%%%%%%%%%%%%%%%%%%%%%%%
\summary

Summarize your thesis.

%\end{spacing}

%=======================================================================
% References
%=======================================================================
\bibliographystyle{ieeetr}
\bibliography{references}

% DO NOT DELETE: 초록 머리말의 전체 쪽수("XXp")는 이 라벨이 있는 쪽(참고문헌 마지막 쪽)으로
% 계산됩니다. 지우면 쪽수가 "??p"로 표시됩니다.
\label{paperlastpagelabel}

%=======================================================================
% Appendix (필요 없으면 아래 3줄을 주석 처리)
%=======================================================================
\begin{spacing}{2.0}
%%%%%%%%%%%%%%%%%%%%%%%%%%%%%%%%%%%%%%%%%%%%%%%%
% Appendix
% \appendix 이후의 \chapter는 "Appendix A", "Appendix B" ... 로 번호가 붙습니다.
%%%%%%%%%%%%%%%%%%%%%%%%%%%%%%%%%%%%%%%%%%%%%%%%
\appendix

\chapter{Supplementary Tables}
%----------------------------------------------------------
% [부록 예시 1] longtblr: 여러 쪽에 걸쳐 이어지는 긴 표
% 표 아래에 작은 글씨로 설명을 덧붙이는 방법도 함께 보여 줍니다.
%----------------------------------------------------------
\begin{longtblr}[
    caption = {Configuration of hyperparameters (example of a \texttt{longtblr} table).},
    label = {tab:hyperparameters},
]{
    width = \linewidth,
    colspec = { Q[c,m] X[l,m] Q[c,m] Q[c,m] },
    row{1} = {font=\bfseries, halign=c},
    rowhead = 1,
}
\toprule
Model   & Hyperparameter & Search space & Selected value \\
\midrule
Model 1 & Learning rate  & $10^{-5}$ to $10^{-2}$ & $10^{-3}$ \\
        & Batch size     & 16, 32, 64   & 32 \\
\addlinespace
Model 2 & Number of trees & 100 to 1000 & 500 \\
        & Maximum depth   & 3 to 10     & 6 \\
\bottomrule
\end{longtblr}

{
    \small
    \setlength{\parindent}{0pt}
    \vspace{1ex}
    \textbf{Note:} Additional explanations for \Cref{tab:hyperparameters} can be written here in a smaller font.
}


\clearpage

\chapter{Supplementary Figures}
%----------------------------------------------------------
% [부록 예시 2] \captionof: figure 환경 없이(떠다니지 않고) 그 자리에 그림 배치
% 부록에 그림을 한 쪽씩 순서대로 넣을 때 유용합니다.
%----------------------------------------------------------
\captionsetup{justification=centering, hypcap=false}

\captionof{figure}{Supplementary figure 1 (example of \textbackslash\texttt{captionof}).}
\label{fig:supp_1}
\includegraphics[width=\linewidth]{example-image-a}
\\[1ex]
\textmd{\small Description of the supplementary figure can be written here.}

\newpage

\captionof{figure}{Supplementary figure 2.}
\label{fig:supp_2}
\includegraphics[width=\linewidth]{example-image-b}
\\[1ex]
\textmd{\small Description of the supplementary figure can be written here.}


\end{spacing}

%=======================================================================
% Acknowledgements -> misc/acknowledgements.tex
%=======================================================================
\begin{acknowledgements}
\begin{spacing}{2.0}
%%%%%%%%%%%%%%%%%%%%%%%%%%%%%%%%%%%%%%
% Acknowledgements  [EDIT]
%%%%%%%%%%%%%%%%%%%%%%%%%%%%%%%%%%%%%%
Write your acknowledgements here. You may thank your advisor, committee members, colleagues, funding agencies, and family.

% (선택) 고지 문단 예시 - 해당되는 경우에만 사용하고, 학과·학교 지침을 따르세요.
% 1) 학위논문 내용 일부가 학술지에 게재되었거나 심사 중인 경우
% Part of this thesis (Chapter X) has been published in / is currently under review in [Journal Name].
% 2) 생성형 AI 도구를 문장 교정 등에 사용한 경우
% A generative AI tool ([Tool name]) was used to improve the readability of the manuscript; all content was reviewed and verified by the author.

\end{spacing}
\end{acknowledgements}

\end{document}
